% =====================================================================
\chapter{Autovalores, autovectores y teorema espectral}
\label{tema:3}
% =====================================================================

Un autovector de una matriz cuadrada es una dirección que la matriz no gira, solo estira o encoge. Este tema define autovalores y autovectores, estudia la diagonalización y demuestra el resultado central para las aplicaciones: las matrices reales simétricas tienen autovalores reales y una base ortonormal de autovectores. Se sigue a Strang~\cite[§6.1--6.3]{strang-ila}, las clases~21, 22 y~25 de 18.06~\cite{mit1806}, las clases~4 y~5 de 18.065~\cite{mit18065} y los §4.2 y~§4.4 de~\cite{mml}.

\section{Definición y polinomio característico}

\begin{definicion}[Autovalor y autovector]
Sea $A\in\R^{n\times n}$. Un número $\lambda\in\mathbb{C}$ es un \emph{autovalor} de $A$ si existe $v\ne 0$ tal que
\[
Av=\lambda v.
\]
Todo $v\ne 0$ que cumple la igualdad es un \emph{autovector} asociado a $\lambda$. El conjunto $\Nuc(A-\lambda I)$, formado por los autovectores y el vector nulo, es el \emph{espacio propio} de $\lambda$.
\end{definicion}

\begin{proposicion}[Polinomio característico]\label{prop:polcar}
$\lambda$ es autovalor de $A$ si y solo si $\det(A-\lambda I)=0$. La función $p_A(\lambda)=\det(A-\lambda I)$ es un polinomio de grado $n$ en $\lambda$, llamado \emph{polinomio característico}.
\end{proposicion}

\begin{proof}
$Av=\lambda v$ con $v\ne 0$ equivale a $(A-\lambda I)v=0$ con $v\ne 0$, es decir, a que $A-\lambda I$ tenga núcleo no trivial. Por el \cref{teo:invertible}, esto ocurre si y solo si $\det(A-\lambda I)=0$. Que $p_A$ es un polinomio de grado $n$ se sigue del desarrollo por cofactores: el único término de grado $n$ procede del producto de la diagonal, $\prod_i(a_{ii}-\lambda)$.
\end{proof}

Por el teorema fundamental del álgebra, $p_A$ tiene $n$ raíces complejas contadas con multiplicidad: $p_A(\lambda)=\prod_{i=1}^{n}(\lambda_i-\lambda)$.

\begin{ejemplo}\label{ej:S}
Sea $S=\begin{pmatrix}2&1\\1&2\end{pmatrix}$. Entonces $p_S(\lambda)=(2-\lambda)^{2}-1=0$, luego $2-\lambda=\pm 1$ y los autovalores son $\lambda_1=3$ y $\lambda_2=1$.
\begin{itemize}
  \item $\lambda_1=3$: $(S-3I)v=0$ da $-v_1+v_2=0$, y un autovector unitario es $q_1=(1,1)/\sqrt{2}$.
  \item $\lambda_2=1$: $(S-I)v=0$ da $v_1+v_2=0$, y un autovector unitario es $q_2=(1,-1)/\sqrt{2}$.
\end{itemize}
Los dos autovectores son ortogonales, $q_1\T q_2=0$. El \cref{teo:simetricas} muestra que no es casualidad.
\end{ejemplo}

\begin{proposicion}[Traza y determinante]\label{prop:traza-det}
Si $\lambda_1,\dots,\lambda_n$ son los autovalores de $A$ contados con multiplicidad, entonces
\[
\det(A)=\prod_{i=1}^{n}\lambda_i,\qquad \traza(A)=\sum_{i=1}^{n}a_{ii}=\sum_{i=1}^{n}\lambda_i.
\]
\end{proposicion}

\begin{proof}
Evaluando $\det(A-\lambda I)=\prod_i(\lambda_i-\lambda)$ en $\lambda=0$ se obtiene la primera igualdad. Para la segunda se compara el coeficiente de $\lambda^{n-1}$ en ambos miembros. En el producto $\prod_i(\lambda_i-\lambda)$ vale $(-1)^{n-1}\sum_i\lambda_i$. En el desarrollo de $\det(A-\lambda I)$, todo término que no sea el producto de la diagonal omite al menos dos factores diagonales y tiene grado $\le n-2$; el producto $\prod_i(a_{ii}-\lambda)$ aporta $(-1)^{n-1}\sum_i a_{ii}$.
\end{proof}

En el \cref{ej:S}, $3+1=4=\traza(S)$ y $3\cdot 1=3=\det(S)$. Estas dos comprobaciones detectan la mayoría de los errores de cálculo.

\begin{table}[h]
\centering\small
\caption{Autovalores de algunas matrices de orden 2.}\label{tab:autovalores}
\begin{tabular}{@{}lcl@{}}
\toprule
Matriz & Autovalores & Interpretación \\
\midrule
$\begin{pmatrix}2&1\\1&2\end{pmatrix}$ & $3$ y $1$ & estira por~3 la diagonal y conserva la antidiagonal \\[6pt]
$\begin{pmatrix}3&1\\0&2\end{pmatrix}$ & $3$ y $2$ & triangular: los autovalores son la diagonal \\[6pt]
$\begin{pmatrix}1&0\\0&0\end{pmatrix}$ & $1$ y $0$ & proyección: conserva el eje $x$ y anula el eje $y$ \\[6pt]
$\begin{pmatrix}0&-1\\1&0\end{pmatrix}$ & $\pm i$ & giro de $90^{\circ}$: ninguna dirección real se conserva \\
\bottomrule
\end{tabular}
\end{table}

\section{Diagonalización}

\begin{proposicion}\label{prop:indep-autovectores}
Autovectores asociados a autovalores distintos son linealmente independientes.
\end{proposicion}

\begin{proof}
Supóngase lo contrario y sea $k$ el menor número de autovectores $v_1,\dots,v_k$, con autovalores distintos $\lambda_1,\dots,\lambda_k$, que admiten una combinación nula $\sum_{i=1}^{k}c_iv_i=0$ con todos los $c_i\ne 0$ (por minimalidad, ninguno puede ser nulo; además $k\ge 2$, porque $v_1\ne 0$). Aplicando $A$ se obtiene $\sum_i c_i\lambda_iv_i=0$, y restando $\lambda_k$ veces la combinación original, $\sum_{i=1}^{k-1}c_i(\lambda_i-\lambda_k)v_i=0$. Los coeficientes $c_i(\lambda_i-\lambda_k)$ son no nulos, lo que da una combinación nula de $k-1$ autovectores y contradice la minimalidad de $k$.
\end{proof}

\begin{teorema}[Diagonalización]\label{teo:diagonalizacion}
$A\in\R^{n\times n}$ admite una factorización $A=X\Lambda X^{-1}$, con $\Lambda$ diagonal y $X$ invertible, si y solo si $A$ tiene $n$ autovectores linealmente independientes. En ese caso, las columnas de $X$ son autovectores, la diagonal de $\Lambda$ contiene los autovalores correspondientes y $A^{k}=X\Lambda^{k}X^{-1}$ para todo $k\ge 1$.
\end{teorema}

\begin{proof}
La igualdad $AX=X\Lambda$ significa, columna a columna, $Ax_i=\lambda_ix_i$. Si los $x_i$ son independientes, $X$ es invertible (\cref{teo:invertible}) y $A=X\Lambda X^{-1}$. Recíprocamente, si $A=X\Lambda X^{-1}$, entonces $AX=X\Lambda$ y las columnas de $X$, independientes por ser $X$ invertible, son autovectores. Finalmente, $A^{k}=X\Lambda X^{-1}X\Lambda X^{-1}\cdots X\Lambda X^{-1}=X\Lambda^{k}X^{-1}$.
\end{proof}

Por la \cref{prop:indep-autovectores}, una matriz con $n$ autovalores distintos es diagonalizable. Con autovalores repetidos puede no serlo: $\begin{pmatrix}1&1\\0&1\end{pmatrix}$ tiene el autovalor doble~1, pero su espacio propio es $\gen\{e_1\}$, de dimensión~1. La diagonalizabilidad no guarda relación con la invertibilidad (\cref{eje:3-mml}).

\section{Matrices simétricas y teorema espectral}

\begin{teorema}\label{teo:simetricas}
Sea $S\in\R^{n\times n}$ simétrica, $S\T=S$. Entonces:
\begin{enumerate}
  \item todos los autovalores de $S$ son reales;
  \item autovectores asociados a autovalores distintos son ortogonales.
\end{enumerate}
\end{teorema}

\begin{proof}
La parte (i) se demuestra en el \cref{eje:3-reales}. Para (ii), sean $Sv_1=\lambda_1v_1$ y $Sv_2=\lambda_2v_2$ con $\lambda_1\ne\lambda_2$. Se calcula $v_1\T Sv_2$ de dos formas:
\[
v_1\T Sv_2=v_1\T(\lambda_2v_2)=\lambda_2\,v_1\T v_2,\qquad
v_1\T Sv_2=(S\T v_1)\T v_2=(Sv_1)\T v_2=\lambda_1\,v_1\T v_2.
\]
Restando, $(\lambda_1-\lambda_2)\,v_1\T v_2=0$ y, como $\lambda_1\ne\lambda_2$, $v_1\T v_2=0$.
\end{proof}

\begin{teorema}[Teorema espectral]\label{teo:espectral}
Toda matriz real simétrica $S\in\R^{n\times n}$ se factoriza como
\[
S=Q\Lambda Q\T=\sum_{i=1}^{n}\lambda_i\,q_iq_i\T,
\]
donde $Q=[q_1\mid\cdots\mid q_n]$ es ortogonal y $\Lambda=\diag(\lambda_1,\dots,\lambda_n)$ es real.
\end{teorema}

\begin{proof}[Demostración en el caso de autovalores distintos]
Por el \cref{teo:simetricas}, los autovalores son reales y los autovectores, que pueden tomarse reales y unitarios, son ortogonales dos a dos. Forman, pues, una base ortonormal, y el \cref{teo:diagonalizacion} con $X=Q$ y $X^{-1}=Q\T$ da $S=Q\Lambda Q\T$. La forma de suma se obtiene multiplicando $Q\Lambda$ por $Q\T$ como columnas por filas. El caso general, con autovalores repetidos, se demuestra mediante la factorización de Schur~\cite[§6.3]{strang-ila}: en cada espacio propio se elige una base ortonormal por Gram--Schmidt.
\end{proof}

Leída de derecha a izquierda, $S=Q\Lambda Q\T$ dice que $S$ expresa el vector en la base de autovectores ($Q\T$), estira cada eje por su autovalor ($\Lambda$) y vuelve a la base original ($Q$). Para el \cref{ej:S},
\[
S=3\,q_1q_1\T+1\,q_2q_2\T=\frac{3}{2}\begin{pmatrix}1&1\\1&1\end{pmatrix}+\frac{1}{2}\begin{pmatrix}1&-1\\-1&1\end{pmatrix}.
\]

\begin{definicion}[Matriz semidefinida positiva]
Una matriz simétrica $S$ es \emph{semidefinida positiva} si $x\T Sx\ge 0$ para todo $x$, y \emph{definida positiva} si $x\T Sx>0$ para todo $x\ne 0$.
\end{definicion}

\begin{proposicion}\label{prop:AtA}
Para toda $A\in\R^{m\times n}$, la matriz $A\T A$ es simétrica y semidefinida positiva, y sus autovalores son mayores o iguales que cero.
\end{proposicion}

\begin{proof}
$(A\T A)\T=A\T A$. Para todo $x$, $x\T A\T Ax=(Ax)\T(Ax)=\norm{Ax}^{2}\ge 0$. Si $A\T Aq=\lambda q$ con $\norm{q}=1$, entonces $\lambda=q\T A\T Aq=\norm{Aq}^{2}\ge 0$.
\end{proof}

\begin{proposicion}[Cociente de Rayleigh]\label{prop:rayleigh}
Sea $S$ simétrica con autovalores $\lambda_1\ge\dots\ge\lambda_n$ y autovectores ortonormales $q_1,\dots,q_n$. Entonces
\[
\max_{\norm{x}=1}x\T Sx=\lambda_1,
\]
y el máximo se alcanza en $x=q_1$.
\end{proposicion}

\begin{proof}
Escríbase $x=\sum_i c_iq_i$ con $\sum_i c_i^{2}=\norm{x}^{2}=1$ (\cref{prop:coord-ortonormal}). Por el \cref{teo:espectral}, $x\T Sx=\sum_i\lambda_ic_i^{2}\le\lambda_1\sum_i c_i^{2}=\lambda_1$. Con $x=q_1$ se obtiene $q_1\T Sq_1=\lambda_1$.
\end{proof}

Este resultado es la base del análisis de componentes principales (\cref{tema:5}): la dirección de máxima varianza de unos datos es el primer autovector de su matriz de covarianza.

\section{El método de la potencia}

Si $S$ es simétrica con $\abs{\lambda_1}>\abs{\lambda_2}\ge\dots\ge\abs{\lambda_n}$, la iteración
\[
x_{k+1}=\frac{Sx_k}{\norm{Sx_k}}
\]
converge, salvo signo, al autovector $q_1$ siempre que $x_0$ tenga componente no nula sobre $q_1$. El error decrece como $\abs{\lambda_2/\lambda_1}^{k}$ (\cref{eje:3-potencia}).

\begin{ejemplo}\label{ej:potencia}
Con la matriz $S$ del \cref{ej:S} y $x_0=(1,0)$: $Sx_0=(2,1)$, $S^{2}x_0=(5,4)$, $S^{3}x_0=(14,13)$, $S^{4}x_0=(41,40)$ y $S^{5}x_0=(122,121)$. Normalizado, $S^{5}x_0/\norm{S^{5}x_0}\approx(0{,}7100;\,0{,}7042)$, muy cerca de $q_1\approx(0{,}7071;\,0{,}7071)$. Como $\abs{\lambda_2/\lambda_1}=1/3$, cada iteración divide el error aproximadamente por~3.
\end{ejemplo}

\begin{observacion}[Conexión con los LLM]
El método de la potencia, y sus variantes por bloques, es la forma habitual de calcular unos pocos autovalores o valores singulares dominantes de matrices grandes~\cite[clase~12]{mit18065}: se usa, por ejemplo, para estimar la norma espectral de una matriz de pesos sin calcular su SVD completa. Las matrices de covarianza de los embeddings son simétricas y semidefinidas positivas (\cref{prop:AtA}), de modo que el \cref{teo:espectral} garantiza que tienen ejes principales ortogonales; esa es la idea del \cref{tema:5}.
\end{observacion}

\begin{lecturas}
\item Strang~\cite{strang-ila}, §6.1 «Introduction to Eigenvalues», §6.2 «Diagonalizing a Matrix» y §6.3 «Symmetric Positive Definite Matrices».
\item MIT 18.06~\cite{mit1806}, clases~21 «Eigenvalues and eigenvectors», 22 «Diagonalization and powers of $A$» y 25 «Symmetric matrices and positive definiteness».
\item MIT 18.065~\cite{mit18065}, clases~4 «Eigenvalues and Eigenvectors» y 5 «Positive Definite and Semidefinite Matrices».
\item Deisenroth, Faisal y Ong~\cite{mml}, §4.2 «Eigenvalues and Eigenvectors» (teorema espectral~4.15) y §4.4 «Eigendecomposition and Diagonalization».
\end{lecturas}

\section*{Ejercicios del tema 3}
\addcontentsline{toc}{section}{Ejercicios del tema 3}

\begin{ejercicio}{1}{eje:3-mml}
Determine, para cada una de las matrices
\[
I_2=\begin{pmatrix}1&0\\0&1\end{pmatrix},\quad P=\begin{pmatrix}1&0\\0&0\end{pmatrix},\quad J=\begin{pmatrix}1&1\\0&1\end{pmatrix},\quad N=\begin{pmatrix}0&1\\0&0\end{pmatrix},
\]
si es diagonalizable y si es invertible. Concluya que ninguna de las dos propiedades implica la otra.
\fuente{adaptado de Deisenroth, Faisal y Ong~\cite{mml}, ejercicio~4.5.}
\end{ejercicio}

\begin{ejercicio}{2}{eje:3-autovec}
Calcule los autovalores y autovectores de $B=\begin{pmatrix}4&1\\2&3\end{pmatrix}$. ¿Son ortogonales los autovectores? Relacione la respuesta con el \cref{teo:simetricas}.
\end{ejercicio}

\begin{ejercicio}{2}{eje:3-u2}
Calcule los autovalores de
\[
A=\begin{pmatrix}2&1&4\\0&3&5\\0&0&-1\end{pmatrix},\qquad
B=\begin{pmatrix}0&0&2\\0&1&0\\8&0&0\end{pmatrix},\qquad
C=\begin{pmatrix}3&3&3\\3&3&3\\3&3&3\end{pmatrix},
\]
y compruebe en cada caso la \cref{prop:traza-det}.
\fuente{adaptado de MIT 18.06SC~\cite{mit1806sc}, sesión «Eigenvalues and Eigenvectors», problema~21.2 (matrices cambiadas).}
\end{ejercicio}

\begin{ejercicio}{2}{eje:3-u3}
Una matriz $B\in\R^{3\times 3}$ tiene autovalores $0$, $2$ y $3$. Calcule, si es posible: (a) $\rango(B)$; (b) $\det(B\T B)$; (c) los autovalores de $(B^{2}+I)^{-1}$. (d) ¿Pueden deducirse los autovalores de $B\T B$?
\fuente{adaptado de MIT 18.06SC~\cite{mit1806sc}, sesión «Eigenvalues and Eigenvectors», problema~21.1 (autovalores cambiados).}
\end{ejercicio}

\begin{ejercicio}{3}{eje:3-reales}
Demuestre que todos los autovalores de una matriz real simétrica son reales. Indicación: si $Sv=\lambda v$ con $v\in\mathbb{C}^{n}$ no nulo, calcule $\bar v\T Sv$ de dos formas.
\end{ejercicio}

\begin{ejercicio}{3}{eje:3-potencia}
Sea $S$ simétrica con autovectores ortonormales $q_1,\dots,q_n$ y autovalores $\abs{\lambda_1}>\abs{\lambda_2}\ge\dots\ge\abs{\lambda_n}$, y sea $x_0=\sum_i c_iq_i$ con $c_1\ne 0$. Demuestre que $S^{k}x_0/\norm{S^{k}x_0}$ tiende a $\pm q_1$ y que la distancia a la recta generada por $q_1$ es $O\bigl(\abs{\lambda_2/\lambda_1}^{k}\bigr)$.
\end{ejercicio}
